% TODO make hw stylesheet
\documentclass[10pt]{article}
\usepackage[letterpaper,top=1in,left=1in,right=1in]{geometry}
\usepackage[dvipsnames,svgnames]{xcolor}
\usepackage[shortlabels]{enumitem}
\usepackage[framemethod=TikZ]{mdframed}
\usepackage{amsmath,amssymb,amsthm}
\usepackage[colorlinks]{hyperref}
\usepackage{multicol}
\usepackage{tabularx}
\usepackage{bbm}

\hypersetup{
    colorlinks=true,
    linkcolor=blue,
    filecolor=magenta,
    urlcolor=magenta,
    pdftitle={MAT 344 HW}
}

\usepackage{mathtools}
\usepackage{thmtools}
\usepackage[textsize=footnotesize]{todonotes}
\usepackage{titling}
\usepackage{cancel}

\reversemarginpar
\mdfdefinestyle{mdbluebox}{roundcorner=10pt,innerbottommargin=9pt,
linecolor=blue,backgroundcolor=TealBlue!5,}
\declaretheoremstyle[headfont=\sffamily\bfseries\color{MidnightBlue},
mdframed={style=mdbluebox},]{thmbluebox}

\mdfdefinestyle{mdbloobox}{frametitlefont=\bfseries,innerbottommargin=8pt,
nobreak=true,backgroundcolor=TealBlue!5,linecolor=blue,}
\declaretheoremstyle[headfont=\bfseries\color{MidnightBlue},
mdframed={style=mdbloobox},headpunct={\\[3pt]},postheadspace=0pt,]{thmbloobox}

\mdfdefinestyle{mdredbox}{frametitlefont=\bfseries,innerbottommargin=8pt,
nobreak=true,backgroundcolor=Salmon!5,linecolor=RawSienna,}
\declaretheoremstyle[headfont=\bfseries\color{RawSienna},
mdframed={style=mdredbox},headpunct={\\[3pt]},postheadspace=0pt,]{thmredbox}

\mdfdefinestyle{mdgreenbox}{linecolor=ForestGreen,backgroundcolor=ForestGreen!5,
linewidth=2pt,rightline=false,leftline=true,topline=false,bottomline=false,}
\declaretheoremstyle[headfont=\bfseries\sffamily\color{ForestGreen!70!black},
mdframed={style=mdgreenbox},headpunct={ --- },]{thmgreenbox}

\mdfdefinestyle{mdorangebox}{linecolor=RedOrange,backgroundcolor=RedOrange!5,
linewidth=2pt,rightline=false,leftline=true,topline=false,bottomline=false,}
\declaretheoremstyle[headfont=\bfseries\sffamily\color{RedOrange!70!black},
mdframed={style=mdorangebox},headpunct={ --- },]{thmorangebox}

\mdfdefinestyle{mdblackbox}{linecolor=black,backgroundcolor=RedViolet!5!gray!5,
linewidth=3pt,rightline=false,leftline=true,topline=false,bottomline=false,}
\declaretheoremstyle[mdframed={style=mdblackbox}]{thmblackbox}

\declaretheoremstyle[spaceabove=6pt,spacebelow=6pt]{boldhead}
\declaretheoremstyle[spaceabove=6pt,spacebelow=6pt,headfont=\normalfont\itshape]{italichead}

\declaretheorem[style=boldhead,name=Problem]{problem}
\declaretheorem[style=boldhead,name=Problem,numbered=no]{problem*}
\declaretheorem[style=boldhead,name=Setup]{setup} % for config geo
\declaretheorem[style=boldhead,name=Example,sibling=problem]{example}
\declaretheorem[style=boldhead,name=Theorem,sibling=problem]{theorem}
\declaretheorem[style=boldhead,name=Corollary,sibling=theorem]{corollary}
\declaretheorem[style=boldhead,name=Proposition,sibling=theorem]{prop}
\declaretheorem[style=boldhead,name=Lemma,numberwithin=problem]{lemma}
\declaretheorem[style=boldhead,name=Theorem,numbered=no]{theorem*}
\declaretheorem[style=boldhead,name=Lemma,numbered=no]{lemma*}
\declaretheorem[style=boldhead,name=Claim,numberwithin=problem]{claim}
\declaretheorem[style=boldhead,name=Claim,numbered=no]{claim*}
\declaretheorem[style=boldhead,name=Remark,numberwithin=problem]{remark}
\declaretheorem[style=boldhead,name=Remark,numbered=no]{remark*}
\declaretheorem[style=boldhead,name=Definition,sibling=theorem]{definition}
\declaretheorem[style=boldhead,name=Definition,numbered=no]{definition*}

\providecommand{\ol}{\overline}
\providecommand{\ceil}[1]{\left\lceil #1 \right\rceil}
\providecommand{\floor}[1]{\left\lfloor #1 \right\rfloor}
\newcommand{\dd}{\mathrm{d}}
\providecommand{\ul}{\underline}
\providecommand{\eps}{\varepsilon}
\providecommand{\half}{\frac{1}{2}}
\providecommand{\CC}{\mathbb C}
\providecommand{\FF}{\mathbb F}
\providecommand{\II}{\mathbb I}
\providecommand{\NN}{\mathbb N}
\providecommand{\QQ}{\mathbb Q}
\providecommand{\RR}{\mathbb R}
\providecommand{\ZZ}{\mathbb Z}
\providecommand{\ind}{\mathbbm 1}
\providecommand{\dg}{^\circ}
\providecommand{\ii}{\item}
\providecommand{\setdiff}{\smallsetminus}
\providecommand{\alert}[1]{{\sffamily\textbf{\textcolor{blue}{#1}}}}
\providecommand{\inv}{^{-1}}
\providecommand{\mb}[1]{\mathbf{#1}}
\newcommand{\what}{\widehat}

\newenvironment{soln}{\begin{proof}[Solution]}{\end{proof}}

\providecommand{\pow}{\operatorname{Pow}}
\providecommand{\vocab}[1]{{\textbf{\textcolor{ForestGreen}{#1}}}}
\providecommand{\lcm}{\operatorname{lcm}}
\providecommand{\abs}[1]{\left\lvert #1\right\rvert}
\providecommand{\smabs}[1]{\lvert #1\rvert}
\providecommand{\innerprod}[1]{\langle\!\langle #1\rangle\!\rangle}
\providecommand{\norm}[1]{\left\| #1\right\|}
\providecommand{\smnorm}[1]{\| #1\|}
\providecommand{\gr}{\operatorname{gr}}
\providecommand{\im}{\operatorname{im}}
\providecommand{\paren}[1]{\left(#1\right)}

\usepackage{mathrsfs}

% place title at top right
\setlength{\droptitle}{-8ex}
\pretitle{\begin{flushright}\Large\bfseries}
\posttitle{\par\end{flushright}}
\preauthor{\begin{flushright}}
\postauthor{\end{flushright}}
\predate{\begin{flushright}}
\postdate{\end{flushright}}

\title{MAT 344 HW09}
\author{\large Aditya Pahuja}
\date{\large Due 04/21}
\begin{document}
\maketitle

\begin{problem}
    [p.200\#1]
    Prove the formula of Gauss:
    \begin{equation*}
	(2\pi)^{\frac{n-1}2}\Gamma(z)
	= n^{z - \half} \Gamma\paren{\frac zn}\Gamma\paren{\frac{z+1}n}
	\cdots\Gamma\paren{\frac{z+n-1}n}.
    \end{equation*}
\end{problem}

\begin{soln}
    Observe that
    \begin{align*}
	\sum_{i=0}^{n-1} \frac{\dd}{\dd z}
	\paren{\frac{\Gamma'(\frac zn + \frac in)}
	{\Gamma(\frac zn + \frac in)}}
	&= \sum_{i=1}^n \sum_{k=1}^\infty
	\frac 1{(\frac zn + \frac in + k)^2}\\
	&= n^2\sum_{i=1}^n \sum_{k=1}^\infty \frac 1{(z + nk + i)^2}
	= n^2\sum_{m=1}^\infty \frac 1{(z + m)^2}
	= n^2\frac{\dd}{\dd z}\paren{\frac{\Gamma'(z)}{\Gamma(z)}}
    \end{align*}
    so
    \begin{equation*}
        \Gamma\paren{\frac zn}\Gamma\paren{\frac{z+1}n}
	\cdots\Gamma\paren{\frac{z+n-1}n} = e^{az + b}\Gamma(z)
    \end{equation*}
    for some constants $a$ and $b$. Taking $z = 1$ and $z = 2$, we get
    \begin{align*}
        \Gamma\paren{\frac 1n}\Gamma\paren{\frac{2}n}
	\cdots\Gamma\paren{\frac{n-1}n}
	&= e^{a + b}\\
        \frac 1n\cdot \Gamma\paren{\frac 1n}\Gamma\paren{\frac{2}n}
	\cdots\Gamma\paren{\frac{n-1}n}
	&= e^{2a + b}
    \end{align*}
    so $e^a = \frac 1n$ and
    \begin{equation*}
	e^b = n
        \cdot \Gamma\paren{\frac 1n}\Gamma\paren{\frac{2}n}
	\cdots\Gamma\paren{\frac{n-1}n}
	= n\sqrt{\prod_{k=1}^{n-1}
	\Gamma\paren{\frac kn}\Gamma\paren{1- \frac kn}}
	= \frac{n\pi^{\frac{n-1}2}}{\sqrt{\prod_{k=1}^{n-1}
	\sin(\frac{\pi k}n)}}.
    \end{equation*}
    Writing $\sin(x) = \frac 1{2i}(e^{ix} - e^{-ix})$, we get
    \begin{equation*}
	\prod_{k=1}^{n-1} \sin\paren{\frac{\pi k}n} = \frac 1{(2i)^{n-1}}
	\prod_{k=1}^{n-1} (e^{\pi k i/n} - e^{-\pi k i/n}).
    \end{equation*}
    Letting $\omega = e^{2\pi i/n}$, we have
    \begin{equation*}
	(e^{\pi k i/n} - e^{-\pi ki/n})
	(e^{\pi (n-k) i/n} - e^{-\pi (n-k)i/n})
	= -2\cos(\pi(n - 2k)/n) - 2
	= \omega^k + \omega^{-k} - 2.
    \end{equation*}
    This means that the square of the product of the
    $e^{\pi ki/n} - e^{-\pi ki/n}$ is equal to
    \begin{equation*}
	\prod_{k=1}^{n-1} (\omega^k + \omega^{-k} - 2)
	= \prod_{k=1}^{n-1} \frac{(\omega^k - 1)^2}{\omega^k}
	= \frac{n^2}{\omega^{n(n-1)/2}}.
    \end{equation*}
    Knowing that the product of the sines is a positive real number,
    we can conclude that it is equal to $\frac{n}{2^{n-1}}$, so
    \begin{equation*}
	e^b = n^{\half} (2\pi)^{\frac{n-1}2},
    \end{equation*}
    which proves the desired identity after substituting in
    the values of $e^a$ and $e^b$ and rearranging.
\end{soln}

\begin{problem}
    [p.200\#2]
    Show that
    \begin{equation*}
	\Gamma\paren{\frac 16} = 2^{-\frac 13}\paren{\frac 3\pi}^{\half}
	\Gamma\paren{\frac 13}^2.
    \end{equation*}
\end{problem}

\begin{soln}
    By Legendre's duplication formula applied to $z = \frac 16$
    and the identity for $\Gamma(z)\Gamma(1-z)$,
    \begin{equation*}
	\sqrt\pi\Gamma\paren{\frac 13}
	= 2^{-\frac 23}\Gamma\paren{\frac 16}\Gamma\paren{\frac 23}
	= 2^{-\frac 23}\Gamma\paren{\frac 16}
	\frac{2\pi}{\sqrt 3\cdot \Gamma(\frac 13)},
    \end{equation*}
    which gives us what we want after rearranging.
\end{soln}

\begin{problem}
    [p.200\#3]
    What are the residues of $\Gamma(z)$ at the poles $z = -n$?
\end{problem}

\begin{soln}
    Looking at the product expansion of $\Gamma(z)$
    shows that its poles are all simple.
    For an arbitrary nonnegative integer $n$,
    \begin{equation*}
	\lim_{z\to -n} (z+ n)\Gamma(z)
	= \lim_{z\to -n} (z + n)
	\frac{\Gamma(z + n + 1)}{z(z + 1)\cdots(z + n)}
	= \lim_{z\to -n} \frac{\Gamma(z + n + 1)}{z(z + 1)\cdots(z + n - 1)}
	= \boxed{\frac{(-1)^n}{n!}}.\qedhere
    \end{equation*}
\end{soln}

\begin{problem}
    Show the convergence of the series of $\gamma$, namely
    \begin{equation*}
	\gamma - 1 = \lim_{N\to\infty}\paren{
	    \sum_{n=1}^N \frac 1{n+1} - \log(N + 1)
	}.
    \end{equation*}
\end{problem}

\begin{soln}
    Write the argument of the limit as
    \begin{align*}
	\sum_{n=1}^N \frac 1{n+1} - \log\paren{1 + \frac 1n}
	&= \frac 1{N+1} - \log 2
	+ \sum_{n=2}^N \frac 1n - \log\paren{1 + \frac 1n}\\
	&= \frac 1{N+1} - \log 2
	+ \sum_{n=2}^N\paren{\frac 1n - \frac 1n + \frac 2{n^2}
	- \frac 3{n^3} + \cdots}\\
	&= \frac 1{N+1} - \log 2 +
	\sum_{n=2}^N \paren{\frac 2{n^2} - \frac 3{n^3} +
	\frac 4{n^4} -\cdots}.
    \end{align*}
    The $n$th summand in the sum at the end is nonnegative
    and bounded above by $\frac 2{n^2}$,
    so its limit as $N$ goes to infinity converges
    by comparison with $\sum_{n=2}^\infty \frac 2{n^2} = \frac{\pi^2}3 - 2$,
    hence
    \begin{equation*}
	\lim_{N\to\infty}\paren{
	    \sum_{n=1}^N \frac 1{n+1} - \log(N + 1)
	}
	= \sum_{n=2}^\infty \paren{\frac 2{n^2} - \frac 3{n^3} +
	\frac 4{n^4} -\cdots} - \log 2
    \end{equation*}
    exists.
\end{soln}










\end{document}
